% sections2/sec7_algebras.tex
\section{\texorpdfstring{$R$}{R}-代数的张量积}
\label{sec:algebras}

在前面的章节中，我们多次见到了代数（环）之间的张量积同构，例如 $\C \otimes_\R \mathrm{M}_n(\R) \cong \mathrm{M}_n(\C)$ 作为复向量空间同构.现在我们将证明：两个环的张量积自身可以赋予自然的环乘法结构，从而使这些同构直接在代数（环）层面上成立.

为了将矩阵代数等非交换结构纳入其中，我们允许所讨论的 $R$-代数非交换，但基础标量环 $R$ 始终假定为交换环.

一个 \newterm{r-algebra-def}{$R$-代数}{$R$-algebra} 是一个赋予了 $R$-双线性乘法 $A \times A \to A$ 的 $R$-模 $A$，满足结合律、存在乘法幺元 $1_A$，并且标量乘法与环乘法满足：
\begin{equation}
  r(ab) = (ra)b = a(rb) \quad (\forall r \in R, a, b \in A). \label{eq:7.1}
\end{equation}
由张量积的普遍映射性质，给出双线性乘法 $A \times A \to A$ 等价于给出 $R$-线性映射 $m \colon A \otimes_R A \to A$.

\begin{theorem}[代数张量积的乘法结构]\label{thm:algebra-tensor-mult}
设 $A$ 与 $B$ 为 $R$-代数.在加法群 $A \otimes_R B$ 上存在唯一的环乘法，使其成为具有幺元 $1 \otimes 1$ 的 $R$-代数，且对所有简单张量满足：
\begin{equation}
  (a \otimes b)(a' \otimes b') = aa' \otimes bb'. \label{eq:7.2}
\end{equation}
\end{theorem}

\begin{proof}
利用 $A$ 与 $B$ 自身的乘法线性映射 $m_A \colon A \otimes_R A \to A$ 与 $m_B \colon B \otimes_R B \to B$，形成张量积映射 $m_A \otimes m_B \colon (A \otimes_R A) \otimes_R (B \otimes_R B) \to A \otimes_R B$.结合张量积的结合律与交换律同构：
\[
  (A \otimes_R B) \otimes_R (A \otimes_R B) \cong (A \otimes_R A) \otimes_R (B \otimes_R B) \xrightarrow{m_A \otimes m_B} A \otimes_R B,
\]
它在简单张量上正好给出 $(a \otimes b) \otimes (a' \otimes b') \mapsto aa' \otimes bb'$.由此诱导出的乘法满足结合律且以 $1 \otimes 1$ 为幺元.
\end{proof}

\begin{corollary}\label{cor:comm-algebras-tensor}
若 $A$ 与 $B$ 为交换 $R$-代数，则 $A \otimes_R B$ 亦为交换 $R$-代数.
\end{corollary}

\begin{example}\label{ex:c-tensor-c}
考虑实代数 $\C \otimes_\R \C$.它作为实向量空间是 $4$ 维的，具有基 $\{1 \otimes 1, 1 \otimes i, i \otimes 1, i \otimes i\}$.令 $x = i \otimes 1$ 与 $y = 1 \otimes i$，则 $x^2 = -1, y^2 = -1$ 且 $xy = yx$.由于方程 $T^2 + 1 = 0$ 在其中有四个根 $(\pm x, \pm y)$，该交换环绝不是域！事实上，存在显式代数同构：
\[
  \C \otimes_\R \C \cong \C \times \C, \quad z \otimes w \mapsto (zw, z\bar{w}).
\]
\end{example}

\begin{example}\label{ex:r-tensor-q-r-zero-divisors}
代数 $\R \otimes_\Q \R$ 拥有大量的零因子：
\[
  (\sqrt{2} \otimes 1 + 1 \otimes \sqrt{2})(\sqrt{2} \otimes 1 - 1 \otimes \sqrt{2}) = 2 \otimes 1 - 1 \otimes 2 = 0,
\]
而左右两因子均非零.
\end{example}

\begin{example}\label{ex:poly-tensor-algebra}
作为 $R$-代数，有自然同构
\[
  R[X] \otimes_R R[Y] \cong R[X, Y], \quad R[X]^{\otimes k} \cong R[X_1, \dots, X_k].
\]
\end{example}

\begin{example}\label{ex:rational-function-field-tensor}
对于域 $F$ 上纯超越扩域的张量积：$F(x) \otimes_F F(y)$ \textbf{绝不是域}，因而绝不同构于 $F(x, y)$！事实上，它同构于分母可分解为 $b(x)d(y)$ 形式的有理函数子环（即局部化环）.元素 $1 \otimes 1 - x \otimes y$ 在其中不可逆（因为 $1/(1-xy)$ 无法分拆分母）.
\end{example}

\begin{theorem}[代数张量积的普遍映射性质]\label{thm:algebra-tensor-ump}
设 $A$ 与 $B$ 为 $R$-代数.对于任意 $R$-代数 $C$ 以及任意满足像元对易（即 $\varphi(a)\psi(b) = \psi(b)\varphi(a)$）的 $R$-代数同态 $\varphi \colon A \to C$ 与 $\psi \colon B \to C$，存在唯一的 $R$-代数同态
\[
  \Phi \colon A \otimes_R B \to C,
\]
满足 $\Phi(a \otimes b) = \varphi(a)\psi(b)$.
\end{theorem}

\begin{theorem}[多项式代数商环的基变换]\label{thm:poly-quotient-algebra-base-ext}
设 $f \colon R \to S$ 为交换环同态，$h(X) \in R[X]$.则作为 $S$-代数有
\begin{equation}
  S \otimes_R (R[X]/(h(X))) \cong S[X]/(h^f(X)),
\end{equation}
其中 $h^f(X)$ 表示将 $h(X)$ 的系数通过 $f$ 映到 $S$ 中所得的多项式.
\end{theorem}

\begin{example}\label{ex:field-extension-tensors}
利用定理~\ref{thm:poly-quotient-algebra-base-ext}，可以极其简便地计算数域扩张的张量积：
\begin{enumerate}
  \item $\R \otimes_\Q \Q(\sqrt{2}) \cong \R[X]/(X^2 - 2) \cong \R \times \R$；
  \item $\R \otimes_\Q \Q(\sqrt[3]{2}) \cong \R[X]/(X^3 - 2) \cong \R \times \C$；
  \item $\Q(\sqrt{2}) \otimes_\Q \Q(\sqrt{3}) \cong \Q(\sqrt{2})[X]/(X^2 - 3) \cong \Q(\sqrt{2}, \sqrt{3})$（这是一个域！）.
\end{enumerate}
\end{example}

\begin{theorem}\label{thm:intermediate-field-tensor}
设 $L/K$ 与 $F/K$ 为域扩张，且 $K \subset E \subset L$ 为中间域.则有典范环同构
\[
  F \otimes_K L \cong (F \otimes_K E) \otimes_E L.
\]
\end{theorem}
